

\section{Results and Discussion}

In this section, we conduct a series of experiments to evaluate the proposed approach for functional module identification in PPI networks. We first investigate the effects of different loss functions and GNN architectures to identify the most suitable model configuration. Based on the results of this comparison, we further examine the aggregation mechanisms of JKNet and conduct an ablation study to evaluate the effect of different activation functions on model performance. The experimental analysis is organized into the following stages:

\begin{itemize}
	\item Comparing the performance of different loss functions and GNN architectures
	\item Analyzing the effect of different JKNet aggregation mechanisms
	\item Comparing the performance of different activation functions
\end{itemize}

These experiments are conducted on three PPI datasets: Collins, Krogan-Core, and DIP, which are used for model architecture selection. BioGRID and Krogan-Extended are excluded from this stage and are subsequently used in the final comparison with existing methods. The improvements observed in the final results on these additional datasets provide further evidence supporting the effectiveness of the selected architecture and its suitability for functional module identification. In all experiments, binary GO features are used as input features from MF and BP subgraphs, and a threshold of $0.3$ is applied to the affiliation matrix to identify functional modules. The models are trained for 2,000 epochs on the Collins and Krogan-Core datasets and 5,000 epochs on the DIP dataset to accommodate the differences in dataset size.

Following these experiments, we evaluate the proposed threshold aggregation strategy for extracting functional modules from the affiliation matrix. Finally, we compare the performance of the selected model with existing methods from the literature to assess its effectiveness relative to previously proposed approaches.

\begin{table*}[h]
	\centering
	\caption{Performance of the GNN architectures and loss functions across the Collins, Krogan-Core, and DIP datasets}
	\vspace{2mm}
	\label{table:loss_arch_results}
	\footnotesize
	\vspace{2mm}
	\begin{tabular}{c| c c c | c c c c c c}
		\hline
		\textbf{Dataset} & \textbf{GNN} & \textbf{Loss} & $\boldsymbol{\lambda}$
		& \textbf{P} & \textbf{R} & \textbf{F1} & \textbf{Acc} & \textbf{NCP} & \textbf{NCB}
		\\
		\hline
		
		\multirow{8}{*}{\textbf{Collins}}
		& \multirow{4}{*}{GCN}
		& NOCD & N/A
		& 0.891 & 0.337 & 0.489 & 0.252 & 49 & 187
		\\
		
		& & W-NOCD & N/A
		& 0.687 & \textbf{0.695} & 0.691 & \textbf{0.326} & 167 & \textbf{385}
		\\
		
		& & HO-NOCD & 0.5
		& 0.818 & 0.316 & 0.456 & 0.254 & 54 & 175
		\\
		
		& & HO-NOCD & 1
		& 0.662 & 0.212 & 0.322 & 0.245 & 49 & 118
		\\
		\cline{2-10}
		
		& \multirow{4}{*}{GAT}
		& NOCD & N/A
		& 0.825 & 0.522 & 0.640 & 0.295 & 156 & 289
		\\
		
		& & W-NOCD & N/A
		& 0.879 & 0.630 & \textbf{0.734} & 0.292 & \textbf{240} & 349
		\\
		
		& & HO-NOCD & 0.5
		& 0.880 & 0.525 & 0.658 & 0.282 & 161 & 291
		\\
		
		& & HO-NOCD & 1
		& \textbf{0.916} & 0.381 & 0.538 & 0.276 & 186 & 211
		\\
		
		\hline
		
		\multirow{8}{*}{\textbf{Krogan-Core}}
		& \multirow{4}{*}{GCN}
		& NOCD & N/A
		& 0.435 & 0.555 & 0.488 & 0.290 & 141 & 333
		\\
		
		& & W-NOCD & N/A
		& 0.478 & 0.603 & 0.533 & \textbf{0.294} & 153 & 362
		\\
		
		& & HO-NOCD & 0.5
		& 0.504 & 0.428 & 0.463 & 0.265 & 123 & 257
		\\
		
		& & HO-NOCD & 1
		& 0.632 & 0.373 & 0.469 & 0.251 & 117 & 224
		\\
		\cline{2-10}
		
		& \multirow{4}{*}{GAT}
		& NOCD & N/A
		& 0.598 & 0.543 & 0.569 & 0.257 & 238 & 326
		\\
		
		& & W-NOCD & N/A
		& \textbf{0.655} & \textbf{0.605} & \textbf{0.629} & 0.261 & \textbf{266} & \textbf{363}
		\\
		
		& & HO-NOCD & 0.5
		& 0.514 & 0.360 & 0.423 & 0.243 & 187 & 216
		\\
		
		& & HO-NOCD & 1
		& 0.308 & 0.220 & 0.257 & 0.205 & 122 & 132
		\\
		
		\hline
		
		\multirow{8}{*}{\textbf{DIP}}
		& \multirow{4}{*}{GCN}
		& NOCD & N/A
		& 0.285 & 0.443 & 0.347 & \textbf{0.278} & 122 & 343
		\\
		
		& & W-NOCD & N/A
		& \textbf{0.528} & 0.045 & 0.485 & 0.257 & 134 & 348
		\\
		
		& & HO-NOCD & 0.5
		& 0.304 & 0.163 & 0.212 & 0.246 & 56 & 126
		\\
		
		& & HO-NOCD & 1
		& 0.233 & 0.134 & 0.170 & 0.227 & 45 & 104
		\\
		\cline{2-10}
		
		& \multirow{4}{*}{GAT}
		& NOCD & N/A
		& 0.314 & 0.413 & 0.357 & 0.242 & 153 & 320
		\\
		
		& & W-NOCD & N/A
		& 0.523 & \textbf{0.506} & \textbf{0.515} & 0.242 & \textbf{212} & \textbf{392}
		\\
		
		& & HO-NOCD & 0.5
		& 0.164 & 0.070 & 0.098 & 0.178 & 22 & 54
		\\
		
		& & HO-NOCD & 1
		& 0.030 & 0.031 & 0.030 & 0.173 & 11 & 24
		\\
		
		\hline
	\end{tabular}
\end{table*}


\subsection{Comparing the performance of loss functions and graph neural network architecture}

In this experiment, we compare the proposed loss functions, W-NOCD and HO-NOCD, with the original NOCD loss using GCN and GAT architectures. The evaluation is conducted on the Collins, Krogan-Core, and DIP datasets, selected to represent different PPI network sizes and characteristics. Binary GO annotations derived from the BP and MF subgraphs are used as input features, and a fixed threshold of $0.3$ is applied to the learned protein–community membership matrix to identify functional modules. The results are presented in Table~\ref{table:loss_arch_results}.

Overall, W-NOCD combined with GAT achieves the highest F1-score across all evaluated datasets. W-NOCD also consistently improves F1-score over the original NOCD loss when used with GCN, highlighting the benefits of incorporating interaction weights into the loss function. In contrast, HO-NOCD does not consistently outperform the original loss, and increasing $\lambda$ generally leads to poorer performance. Although some configurations yield modest improvements, the higher-order formulation provides less consistent benefits than W-NOCD under the settings considered.

The precision and recall results further demonstrate the favorable balance achieved by W-NOCD, particularly when combined with GAT. Although accuracy varies across configurations, the consistent superiority in F1-score supports selecting GAT with W-NOCD as the final model. To further investigate the choice of GNN architecture, we also evaluate GATv2 and SuperGAT using binary GO annotations. As reported in Appendix~\ref{}, neither variant outperforms the standard GAT, further supporting its selection for functional module identification.


\begin{table*}[t]
	\centering
	\caption{Comparison of the GAT baseline and JKNet aggregation mechanisms using binary GO features and the W-NOCD loss}
	\label{table:jknet_binary}
	\footnotesize
	\vspace{2mm}
	\begin{tabular}{c|c|c|c|c|c|c|c}
		\hline
		\textbf{Dataset} & \textbf{Architecture} & \textbf{Precision} &
		\textbf{Recall} & \textbf{F1} & \textbf{Accuracy} &
		\textbf{NCP} & \textbf{NCB} \\
		\hline
		
		\multirow{4}{*}{Collins}
		& GAT (baseline) & 0.879 & 0.630 & \textbf{0.734} & 0.292 & 240 & 349 \\
		& JKNet-MaxPooling & 0.849 & 0.511 & 0.638 & 0.271 & 163 & 283 \\
		& JKNet-Concatenation & 0.795 & \textbf{0.671} & 0.728 & \textbf{0.294} & \textbf{267} & \textbf{372} \\
		& JKNet-LSTM Attention & \textbf{0.898} & 0.509 & 0.650 & 0.268 & 193 & 282 \\
		\hline
		
		\multirow{4}{*}{Krogan-Core}
		& GAT (baseline) & \textbf{0.655} & \textbf{0.605} & \textbf{0.629} & 0.261 & \textbf{266} & \textbf{363} \\
		& JKNet-MaxPooling & 0.534 & 0.572 & 0.552 & 0.270 & 211 & 343 \\
		& JKNet-Concatenation & 0.542 & 0.572 & 0.557 & \textbf{0.269} & 217 & 343 \\
		& JKNet-LSTM Attention & 0.526 & 0.595 & 0.558 & 0.268 & 245 & 357 \\
		\hline
		
		\multirow{4}{*}{DIP}
		& GAT (baseline) & 0.523 & 0.506 & 0.515 & 0.242 & 212 & 392 \\
		& JKNet-MaxPooling & 0.517 & 0.541 & 0.529 & \textbf{0.265} & 194 & 419 \\
		& JKNet-Concatenation & \textbf{0.530} & 0.545 & \textbf{0.537} & 0.254 & 204 & 422 \\
		& JKNet-LSTM Attention & 0.498 & \textbf{0.579} & 0.535 & 0.255 & \textbf{216} & \textbf{448} \\
		\hline
	\end{tabular}
\end{table*}



\subsection{Analyzing the JKNet Mechanism}

To investigate the effect of layer aggregation on functional module identification, we evaluate MaxPooling, concatenation, and LSTM-based attention using binary GO features, the W-NOCD loss, and GAT as the base GNN. Experiments are conducted on the Collins, Krogan-Core, and DIP datasets, with results reported in Table~\ref{table:jknet_binary}.

The results show that concatenation achieves the highest F1-score on Collins and DIP, while LSTM-based attention performs best on Krogan-Core by a small margin. Although the best-performing strategy varies across datasets, concatenation generally provides a favorable balance between precision and recall. MaxPooling does not achieve the highest F1-score on any dataset, while the differences among aggregation strategies are relatively small on Krogan-Core and DIP.

Overall, concatenation is a competitive aggregation mechanism, particularly on Collins, where its advantage is more pronounced. However, none of the JKNet strategies consistently outperforms the standard GAT architecture evaluated in the preceding experiments. This suggests that aggregating representations from multiple layers does not necessarily improve functional module identification in the considered settings.


\begin{table*}[t]
	\centering
	\caption{Performance of different activation functions using binary GO features across three PPI datasets}
	\label{table:activation_binary_go}
	\footnotesize
	\vspace{2mm}
		\begin{tabular}{c|c|c|c|c|c|c|c}
			\hline
			\textbf{Dataset} &
			\textbf{Activation} &
			\textbf{Precision} &
			\textbf{Recall} &
			\textbf{F1} &
			\textbf{Accuracy} &
			\textbf{NCP} &
			\textbf{NCB} \\
			\hline
			\multirow{3}{*}{Collins}
			& ELU  & \textbf{0.922} & 0.403 & 0.560 & 0.240 & 71  & 223 \\
			& GELU & 0.852 & 0.455 & 0.593 & 0.270 & 98  & 252 \\
			& ReLU & 0.889 & \textbf{0.563} & \textbf{0.689} & \textbf{0.286} & \textbf{112} & \textbf{312} \\
			\hline
			\multirow{3}{*}{Krogan-Core}
			& ELU  & \textbf{0.706} & 0.468 & 0.563 & 0.225 & 137 & 281 \\
			& GELU & 0.675 & 0.560 & \textbf{0.612} & 0.254 & \textbf{237} & 336 \\
			& ReLU & 0.632 & \textbf{0.590} & 0.610 & \textbf{0.263} & 208 & \textbf{354} \\
			\hline
			\multirow{3}{*}{DIP}
			& ELU  & \textbf{0.667} & 0.283 & 0.397 & 0.187 & 102 & 219 \\
			& GELU & 0.605 & 0.479 & 0.535 & 0.233 & \textbf{233} & 371 \\
			& ReLU & 0.614 & \textbf{0.522} & \textbf{0.564} & \textbf{0.240} & 199 & \textbf{404} \\
			\hline
		\end{tabular}
\end{table*}

\subsection{Comparing the Activation Functions}

In this experiment, we investigate the effect of ReLU, ELU, and GELU activation functions on functional module identification using binary GO annotations and GO2Vec embeddings across five PPI datasets. The results in Table~\ref{table:activation_function_comparison} illustrate how activation function performance varies with the input representation and dataset characteristics.

For binary GO features, ReLU achieves the highest F1-score on most datasets, while GELU performs best on Krogan-Core and Krogan-Extended. Similarly, ReLU generally provides higher recall, whereas ELU tends to favor precision at the expense of recall. When GO2Vec embeddings are used, ReLU achieves the highest F1-score on four datasets and consistently yields the highest recall, while ELU performs best on Collins. These findings indicate that ReLU offers a more favorable balance between precision and recall across different datasets and feature representations.

Overall, binary GO features generally outperform GO2Vec embeddings in terms of F1-score, particularly when combined with ReLU. Although ELU and GELU occasionally achieve higher precision or accuracy, these improvements do not consistently translate into better overall performance. Therefore, ReLU is selected for subsequent experiments due to its consistently strong F1-score and recall across the evaluated datasets.





% appendix table:
\begin{table*}[t]
	\centering
	\caption{Performance of different activation functions using GO2Vec features across three PPI datasets}
	\label{table:activation_go2vec}
	\footnotesize
	\vspace{2mm}
	\resizebox{\textwidth}{!}{
		\begin{tabular}{c|c|c|c|c|c|c|c}
			\hline
			\textbf{Dataset} &
			\textbf{Activation} &
			\textbf{Precision} &
			\textbf{Recall} &
			\textbf{F1} &
			\textbf{Accuracy} &
			\textbf{NCP} &
			\textbf{NCB} \\
			\hline
			\multirow{3}{*}{Collins}
			& ELU  & 0.882 & \textbf{0.435} & \textbf{0.583} & 0.257 & \textbf{60} & \textbf{241} \\
			& GELU & 0.892 & 0.366 & 0.520 & \textbf{0.286} & 58  & 203 \\
			& ReLU & \textbf{0.903} & 0.339 & 0.493 & 0.249 & 56  & 188 \\
			\hline
			\multirow{3}{*}{Krogan-Core}
			& ELU  & \textbf{0.902} & 0.225 & 0.360 & 0.188 & 37  & 135 \\
			& GELU & 0.792 & 0.348 & 0.484 & 0.236 & 76  & 209 \\
			& ReLU & 0.658 & \textbf{0.535} & \textbf{0.590} & \textbf{0.264} & \textbf{158} & \textbf{321} \\
			\hline
			\multirow{3}{*}{DIP}
			& ELU  & 0.714 & 0.071 & 0.129 & 0.116 & 15  & 55 \\
			& GELU & \textbf{0.782} & 0.211 & 0.332 & 0.177 & 61  & 163 \\
			& ReLU & 0.674 & \textbf{0.380} & \textbf{0.486} & \textbf{0.220} & \textbf{116} & \textbf{294} \\
			\hline
		\end{tabular}
	}
\end{table*}
